\subsection{Linear mappings with values in a space of continuous functions defined by a kernel}

Here we take K = C. Let X and Y be locally compact topological spaces. Let \(\mu\) be a complex measure on X and k a mapping from X × Y into C possessing the following two properties:

\begin{enumerate}
\def\labelenumi{(\roman{enumi})}
\item
  For every \(y \in Y\) , the function \(k_{y} \colon x \mapsto k(x, y)\) of X into C is \(\mu\) -integrable;
\item
  The map \(y \mapsto k_{y}\) from Y into \(\mathcal{L}_{\mathbf{C}}^{1}(\mathrm{X}, \mu)\) is continuous.
\end{enumerate}

Equip the space \(\mathcal{C}(\mathrm{Y})\) of complex-valued continuous functions on Y with the topology of compact convergence. For \(f\) in \(\mathcal{L}_{\mathbf{C}}^{\infty}(\mathrm{X},\mu)\) , one defines \(\widetilde{k}(f)\colon \mathrm{Y}\to \mathbf{C}\) by

\[
\widetilde {k} (f) (y) = \int_ {\mathrm{X}} k (x, y) f (x) d \mu (x) = \int_ {\mathrm{X}} k _ {y} f d \mu .\tag{1}
\]

The map \(h \mapsto \int_{X} hfd\mu\) is a continuous linear form on \(\mathcal{L}_{\mathbf{C}}^{1}(X,\mu)\) , so condition (ii) shows that the function \(\widetilde{k}(f)\) is continuous.

PROPOSITION 2. — The map \(\widetilde{k}\) from \(\mathcal{L}_{\mathbf{C}}^{\infty}(\mathrm{X},\mu)\) into \(\mathcal{C}(\mathrm{Y})\) is compact.

For any \(y\) and \(y'\) in Y, we have

\[
| \widetilde {k} (f) (y) | \leqslant \| k _ {y} \| _ {1} \| f \| _ {\infty}, \quad | \widetilde {k} (f) (y) - \widetilde {k} (f) (y ^ {\prime}) | \leqslant \| k _ {y} - k _ {y ^ {\prime}} \| _ {1} \| f \| _ {\infty}.
\]

By Ascoli's theorem (TG, X, p. 18, cor. 2), it follows that the image under \(\widetilde{k}\) of the unit ball of \(\mathcal{L}_{\mathbf{C}}^{\infty}(\mathrm{X};\mu)\) is relatively compact in \(\mathcal{C}(\mathrm{Y})\) .

COROLLARY. --- Let X be a compact topological space, Y a locally compact topological space, \(\mu\) a complex measure on X and k a continuous map from X × Y into C. Then formula (1) defines a compact linear map, still denoted \(\widetilde{k}\) , from \(\mathcal{C}(X)\) endowed with the topology of uniform convergence into \(\mathcal{C}(Y)\) endowed with the topology of compact convergence.

Let us verify conditions (i) and (ii) for \(k\) . For every \(y \in \mathrm{Y}\) , the map \(k_y\) is continuous on the compact space X, hence \(\mu\) -integrable. Moreover, the map \(y \mapsto k_y\) is continuous from Y into \(\mathcal{C}(\mathrm{X})\) since X is compact (TG, X, p. 28, th. 3). Since \(\| k_y - k_{y'} \|_1 \leqslant \| \mu \| \| k_y - k_{y'} \|\) for all \((y, y') \in \mathrm{Y}^2\) this mapping from Y into \(\mathcal{L}_{\mathbf{C}}^1(\mathrm{X}, \mu)\) is continuous. By prop. 2, the mapping \(\widetilde{k}\) from \(\mathcal{L}_{\mathbf{C}}^\infty(\mathrm{X}, \mu)\) into \(\mathcal{C}(\mathrm{Y})\) is compact. By composition with the canonical mapping \(\mathcal{C}(\mathrm{X}) \to \mathcal{L}_{\mathbf{C}}^\infty(\mathrm{X}, \mu)\) , one obtains a compact linear mapping from \(\mathcal{C}(\mathrm{X})\) into \(\mathcal{C}(\mathrm{Y})\) .
% semantic-fragments: legacy-input-seam
\relax{}
